\section{The finite parameter cover and the geometric conclusion}
\label{sec:direct-parameter-cover}

The regional estimates bound the displacement on the whole interval
of admissible central heights. It remains to show that their parameter
regions cover a full neighborhood of the organizing value, with common
sections and compatible constant choices. The resulting zero bound is
the maximum of the regional bounds. Geometric capture then gives the
upper bound for periodic orbits. The table below summarizes the regions;
the proof fixes their thresholds and checks their overlaps and boundary
cases.

\begin{table}[H]
\centering
\small
\caption{Parameter regions, the comparison that bounds the zeros of
\(\Delta_p\) on the whole height interval \(\cH_p\), and the resulting
bound. After the first row, \(q_1>0\). The constants
\(L,\kappa_{\rm p},m_{\rm cov},A_0,a_{\rm w},\kappa_{\rm w},r_{\rm d}\) are
chosen in the proof below, in \cref{sec:direct-parameter-cover}, which
also checks the overlaps and the boundary axes.}
\label{tab:direct-cover-guide}
\label{tab:direct-parameter-cover}
\begin{tabular}{@{}>{\raggedright\arraybackslash}p{\dimexpr.31\textwidth-10pt}
                >{\raggedright\arraybackslash}p{\dimexpr.24\textwidth-10pt}
                >{\raggedright\arraybackslash}p{\dimexpr.27\textwidth-10pt}
                >{\raggedright\arraybackslash}p{\dimexpr.1\textwidth-10pt}
                >{\raggedleft\arraybackslash}p{\dimexpr.08\textwidth}@{}}
\hline
Parameter region & Obstruction & Comparison & Section & Zero bound\\
\hline
\(q_1\le0\) & Vanishing radial flux at the central point
  & Signed paired contraction & \ref{sec:paired-variation} & 1\\
\(|q_0|\le Lq_1^2\), arbitrary \(q_2\) & Contact and low flux
  & Scattering ratio and endpoint bounds
  & \ref{sec:direct-bounded-contact} & 2\\
\(q_0>0\), \(q_1/\sqrt{q_0}\le\kappa_{\rm p}\) & Interior saddle
  & Dulac slopes and exterior contraction
  & \ref{sec:direct-positive-offset} & 2\\
\(q_0<0\), \(|q_0|/H^3\ge m_{\rm cov}\) & Merging boundary roots
  & Drift-weighted curvature; early central heights
  & \ref{sec:direct-primary-negative} & 3\\
\(q_0<0\), \(|q_2|/\sqrt{q_1}\le A_0\), \(|q_0|/q_1^{3/2}\le a_{\rm w}\),
\(q_1^2/|q_0|\le\kappa_{\rm w}\) & Weak boundary root
  & Weak-root comparison & \ref{sec:direct-negative-displaced} & 2\\
\(q_0<0\), \(q_2>0\), \(r<r_{\rm d}\), outside the contact strip
  & Near-symmetric roots & Curvature or power comparison
  & \ref{sec:direct-positive-quadratic} & 3\\
\(q_0<0\), \(q_2<0\), \(r<r_{\rm d}\), outside the contact strip
  & Unequal exits; long bottleneck
  & Radially weighted rates and exact boundary action
  & \ref{sec:direct-negative-annulus-buffer}--\ref{sec:direct-negative-root-regimes}
  & 3\\
\hline
\end{tabular}
\end{table}


\begin{proposition}[Uniform zero bound]
\label{prop:direct-uniform-displacement}
For each fixed \(C_0>1\), there is a full parameter
neighborhood of \((1,C_0,0,0,0)\) such that, with
\[
 N_{\rm dir}=3,
\]
\(\Delta_p\) has at most \(N_{\rm dir}\)
distinct zeros on the height interval \(\cH_p\).
\end{proposition}

\begin{proof}
\emph{Common sections and the contraction margin.}
First fix a preliminary box in the coordinates \((\eta,z)\) for this
\(C_0\), the common cone of \cref{sec:direct-passage-estimates}, and the
fixed adaptive and annular landing levels. The organizing endpoints satisfy
\(\eta_\pm^2/(2z_*)\to1\) as \(z_*\downarrow0\), by
\cref{lem:exterior-multiplier-new}. Hence their windows can be
chosen with a fixed lower bound \(z_*/\eta_\pm^2\ge c_{\rm end}>0\)
for all sufficiently small sections; choose each landing level
below \(c_{\rm end}/2\). The other landing restrictions use the enlarged
compact drift domains before the later coefficient and height thresholds.

Choose \(\Lambda_*\) with \(\Lambda_0<\Lambda_*<1\), and reserve a
fixed fraction of \(-\log\Lambda_*\) for tail errors.
By \cref{lem:exterior-multiplier-new} with \(\Lambda_+=\Lambda_*\),
we can then fix \(z_*\) small and choose the enlarged exterior
windows \(\widetilde W_\pm\) of that lemma. Thus
\eqref{eq:exterior-uniform-contraction-new} holds with
\(\Lambda_+=\Lambda_*\) in every regional estimate.

The same tail bound must fit the reserved contraction margin in every
region. The half-orbits enter the common cone after a finite scaled
continuation, as in \cref{lem:early-central-heights},
in \cref{sec:direct-bounded-contact,sec:direct-positive-offset}, and
by \eqref{eq:direct-finite-common-cone-entry}. From a common height
\(t\) in that cone, \eqref{eq:direct-tail-first-slope} bounds the
tail contribution to the logarithm of the endpoint derivative ratio by
\(2C_{\rm cone}(\sqrt{z_*}+|q_1|/\sqrt t)\), as in
\cref{lem:boundary-flow-endpoint}. The coefficient of \(\sqrt{z_*}\)
in every such upper bound is therefore at most \(2C_{\rm cone}\),
independently of \(L\), of coefficient margins such as \(\nu\), and
of the regional cuts.

The continuation constants and higher exterior derivatives may depend
on the cuts and on \(z_*\), but enter only the final parameter
thresholds. Choose \(z_*\) also so that
\(2C_{\rm cone}\sqrt{z_*}\) is smaller than the reserved fraction.
This choice is made after the common cone and before all later cuts
and thresholds, and meets every regional smallness requirement on
\(z_*\). If \(q_1\le0\), \cref{prop:signed-contraction-new} gives
at most one zero. Henceforth assume \(q_1>0\).

\emph{Secondary negative offsets.}
For \(q_0<0,\ q_2\ne0\), let \(s,r\) be the scales of
\cref{sec:standard-tools}. We first fix thresholds for these scales;
the contact width will be chosen after the remaining regions have
been compared.
\Cref{prop:direct-positive-quadratic} for \(q_2>0\) and the second
assertion of \cref{prop:direct-negative-weak-root} for \(q_2<0\) give
constants \(\kappa_+\) and \(\kappa_-\), respectively, and thresholds for
\(s,r\) independent of \(L\). Below these thresholds, for every
\(L>\max\{\kappa_+^{-1},\kappa_-^{-1}\}\), \(\Delta_p\) has at most three
zeros outside the strip \(|q_0|\le Lq_1^2\), for either sign of \(q_2\).
Consequently a single secondary threshold
\(0<r_{\rm d}<1/4\) for \(r\) can be fixed for both signs
before the final contact constant \(L\) is selected.

\emph{The primary and linear scales.}
For the remaining negative offsets, including \(q_2=0\), the primary
and linear scales must overlap the secondary region. We choose the
thresholds so that small normalized constant coefficient, outside the
secondary region, forces the normalized linear coefficient to stay
bounded away from zero. Put
\[
 b_{\rm cov}:=r_{\rm d}/2,\qquad A_0:=2/\sqrt{b_{\rm cov}}.
\]
For this fixed \(A_0\), the primary displaced-contact
estimate \cref{prop:direct-negative-displaced} gives
constants \(a_{\rm w},\kappa_{\rm w}>0\).
Choose \(m_{\rm cov}>0\) so that
\begin{equation}
\label{eq:direct-cover-constants}
 2m_{\rm cov}<1,\qquad \sqrt{2m_{\rm cov}}<b_{\rm cov},
             \qquad 2m_{\rm cov}/b_{\rm cov}^{3/2}<a_{\rm w}.
\end{equation}
With \(H\) and \(\rho_1\) of \cref{sec:standard-tools}, put
\(\bar a=|q_0|/H^3\) and \(\bar b=q_1/H^2\).
If \(\bar a\ge m_{\rm cov}\),
\cref{prop:direct-primary-negative} gives at most three zeros.
On the remaining region \(\bar a<m_{\rm cov}\) (and hence also
\(\bar a\le2m_{\rm cov}\)), use the secondary estimates when
\(q_2\ne0\) and \(r<r_{\rm d}\). For the complementary part,
it suffices to consider \(q_2=0\) or
\(r\ge r_{\rm d}/2=b_{\rm cov}\). On this part we claim
\begin{equation}
\label{eq:direct-cover-primary-overlap}
 \bar b\ge b_{\rm cov},\qquad
 |q_2|/\rho_1\le1/\sqrt{b_{\rm cov}}<A_0,\qquad
 |q_0|/\rho_1^3\le2m_{\rm cov}/b_{\rm cov}^{3/2}<a_{\rm w}.
\end{equation}
To prove the overlap, consider which term realizes the maximum
defining \(H\). If \(H=\rho_1\), the first assertion is immediate.
The possibility \(H=|q_0|^{1/3}\) would give
\(\bar a=1\), excluded by \eqref{eq:direct-cover-constants}.
In the remaining case \(H=|q_2|\),
\[
 r=\max\{\sqrt{\bar a},\bar b\}.
\]
Because \(\sqrt{\bar a}<b_{\rm cov}\) and \(r\ge b_{\rm cov}\),
one has \(\bar b\ge b_{\rm cov}\).
The other assertions in
\eqref{eq:direct-cover-primary-overlap} follow from
\(\rho_1=H\sqrt{\bar b}\).

If also \(q_1^2/|q_0|\le\kappa_{\rm w}\), these inequalities put
the parameter in \cref{prop:direct-negative-displaced}, giving at
most two zeros. Within this part of coefficient space, the
complementary case satisfies \(|q_0|/q_1^2<1/\kappa_{\rm w}\)
and lies in bounded contact. The bands
\(m_{\rm cov}\le\bar a\le2m_{\rm cov}\) and
\(r_{\rm d}/2\le r<r_{\rm d}\) provide fixed overlaps of these
coefficient regions.

\emph{Positive offsets and the contact strip.}
If \(q_0>0\), \cref{prop:direct-positive-offset} gives at most two
zeros whenever \(q_1/\sqrt{q_0}\le\kappa_{\rm p}\), for a fixed
\(\kappa_{\rm p}>0\). Its complement satisfies
\(q_0/q_1^2<\kappa_{\rm p}^{-2}\). We can now choose the contact
width to include all the remaining parameters:
\begin{equation}
\label{eq:direct-final-contact-constant}
 L>2\max\{\kappa_{\rm p}^{-2},\kappa_{\rm w}^{-1},
                         \kappa_+^{-1},\kappa_-^{-1},1\}.
\end{equation}
\Cref{prop:direct-bounded-contact} gives at most two zeros on this
strip for sufficiently small parameters.

\Cref{tab:direct-constant-order} records the global order of choices.
Within each regional proof, its fixed coefficient and height cuts precede
the final scale reduction, as specified there.
\begin{table}[ht]
\centering
\small
\caption{Global order of fixed choices. Each row may depend on preceding
rows; the original parameter neighborhood is reduced last.}
\label{tab:direct-constant-order}
\begin{tabular}{@{}>{\raggedright\arraybackslash}p{\dimexpr.44\textwidth-6pt}
                >{\raggedright\arraybackslash}p{\dimexpr.56\textwidth-6pt}@{}}
\hline
Choice & Purpose and permitted dependence\\
\hline
Fixed \(C_0\); preliminary box, common cone and landing level
\(m_{\rm land}\); \(\Lambda_*=\Lambda_+\), then \(z_*\), exterior windows
and capture data
  & Reserve the exterior margin using \(\Lambda_0\), and fix
    \(2C_{\rm cone}\sqrt{z_*}\) below a fixed fraction of it, before any
    coefficient or height threshold. The landing level is small for
    \cref{lem:direct-adaptive-positive} and
    \eqref{eq:direct-no-root-adaptive-rate}, which use only bounds of the
    drift factors.\\
For \(q_2>0\): \(a_+,\kappa_+\), then the compact thresholds. For
\(q_2<0\): the fold width \(\epsilon_{\rm f}\) and
\(\nu=\epsilon_{\rm f}^2/50\), \(a_-,\kappa_-\), then the compact
thresholds. Then \(r_{\rm d}\)
  & Both signs of \(q_2\); the complements of the weak-root regions lie
    in bounded contact strips.\\
\(b_{\rm cov}=r_{\rm d}/2\), \(A_0\), \(a_{\rm w}\), \(\kappa_{\rm w}\),
\(m_{\rm cov}\)
  & Fixed primary--secondary overlaps in
    \eqref{eq:direct-cover-constants}--\eqref{eq:direct-cover-primary-overlap}.\\
Positive-offset threshold \(\kappa_{\rm p}\), then contact width \(L\)
  & \eqref{eq:direct-final-contact-constant} includes the complements of
    all weak-root and saddle regions.\\
Contact height thresholds, compact passing sets and matching height;
final original parameter neighborhood
  & Cut-dependent finite-height errors and exterior derivative bounds
    are absorbed by parameter reduction with \(z_*\) fixed.\\
\hline
\end{tabular}

\end{table}

\emph{Boundary cases and a full neighborhood.}
The regions now have fixed overlaps. We check that they also include
their boundary cases. If \(q_1=0\), signed contraction applies, including
\(q_0=q_2=0\) and the axes \(H=0\) or \(r=0\). If \(q_1>0\)
and \(q_0=0\), the contact row applies.

If \(q_0<0\) and \(q_2=0\), then \(\bar a<m_{\rm cov}\) puts
\(H=\rho_1\), so \(\bar b=1\) and the primary--contact overlap
above applies. Its weak-root condition puts the point in the displaced
row; the complementary condition puts it in the contact row.

For \(q_0<0\), \(\bar a<m_{\rm cov}\) and \(q_2\ne0\), the
cases \(r<r_{\rm d}/2\), \(r_{\rm d}/2\le r<r_{\rm d}\), and
\(r\ge r_{\rm d}\) are covered by the secondary row, the doubled
overlap, and the primary--contact overlap. The primary row includes
\(\bar a=m_{\rm cov}\); the overlap and secondary rows include
\(r=r_{\rm d}/2\), and the overlap includes \(r=r_{\rm d}\).

Equality \(|q_0|=Lq_1^2\) belongs to contact. For \(q_0>0\), the
saddle row and its contact complement are exhaustive. Equality in the
saddle condition is assigned to the saddle row, and both rows allow
\(q_2=0\). Thus every parameter belongs to at least one row of
\cref{tab:direct-parameter-cover}. The maximum of the listed zero
bounds is \(N_{\rm dir}=3\).

Finally, choose the smallest of the finitely many remaining parameter
thresholds, keeping the sections fixed. The coefficient comparison in
Appendix~\ref{app:coordinates} gives a full neighborhood in the original
five parameters on which all these conditions hold. All estimates are
uniform for \((A,C)\) in the fixed neighborhood, independently of
\(H,s,r\) and \(\rho_1\). This proves the proposition.
\end{proof}

\begin{proof}[Proof of \cref{thm:direct-finite-cyclicity}]
Fix \(C_0>1\) and use the fixed sections, exterior windows,
and full parameter neighborhood constructed in the proof of
\cref{prop:direct-uniform-displacement}. Choose the
whole-graphic capture data of
\cref{prop:isolation-capture-new} for these windows, reduced as in
the paragraph before \cref{lem:exterior-multiplier-new} and in
\eqref{eq:small-central-height-new}.
Let \(V\) be the intersection of the parameter neighborhood
with \(V_{\mathrm{cap}}\), and choose
\(\varepsilon\le\rho_{\mathrm{cap}}\); by
\cref{prop:isolation-capture-new}(ii), every periodic orbit within
Hausdorff distance \(\varepsilon\) of \(\GammaI\) is then captured.

Let \(p\in V\), and let \(\gamma\) be a periodic orbit with
\(\distH(\overline\gamma,\GammaI)<\varepsilon\). By
\cref{lem:unique-central-crossing-new}, it has exactly one transverse
central crossing. The interval construction in
\cref{prop:central-displacement-interval-new} assigns it a zero of
\(\Delta_p\), injectively. The uniform zero bound therefore gives
at most \(N_{\rm dir}=3\) such orbits.

It remains to check isolation. The interior-window construction in
\cref{prop:central-displacement-interval-new} places the height of
every captured orbit in \(\operatorname{int}\mathcal H_p\).
The finite zero bound makes its displacement zero isolated in an
open central-height neighborhood. The correspondence in that
proposition preserves isolation, so every periodic orbit counted
above is a limit cycle.
Consequently \(N(V,\varepsilon)\le3\) in
\eqref{eq:cyclicity-definition}, and \(\Cycl(\GammaI)\le3\).
\end{proof}
